
\SGASubsection{4.3}{Connected Affine Groups in Characteristic \(p>0\)}
\SGARefTarget{sga3:xvii:subsection:4:3}

\SGAStatement{\SGAInlineRefTarget{sga3:xvii:proposition:4-3-1}Proposition 4.3.1}{Let \(G\) be a connected affine
algebraic group over a field \(k\) of characteristic \(p>0\).
The following conditions are equivalent:
\begin{enumerate}
\item[(i)] \SGARefTarget{sga3:xvii:proposition:4:3:1:item:i}
  \(G\) is unipotent.
\item[(ii)] \SGARefTarget{sga3:xvii:proposition:4:3:1:item:ii}
  \(G\) has a composition series whose successive quotients are
  isomorphic to \(\boldsymbol\alpha_p\) and \(\mathbb G_a\),
  taken in that order.
\item[(iii)] \SGARefTarget{sga3:xvii:proposition:4:3:1:item:iii}
  \(G\) has a characteristic composition series whose successive
  quotients are isomorphic to \((\boldsymbol\alpha_p)^r\) and
  \((\mathbb G_a)^s\), taken in that order.
\item[(iv)] \SGARefTarget{sga3:xvii:proposition:4:3:1:item:iv}
  \(G_{\overline{k}}\) contains no subgroup isomorphic to
  \(\boldsymbol\mu_p\).
\item[(v)] \SGARefTarget{sga3:xvii:proposition:4:3:1:item:v}
  \(\mathfrak g=\operatorname{Lie}G\) is unipotent (\SGARefLink{sga3:xvii:definition:3:6}{3.6}).
\item[(vi)] \SGARefTarget{sga3:xvii:proposition:4:3:1:item:vi}
  \(\mathfrak g\) is nilpotent, and the reductive part of the
  center of \(\mathfrak g\) (\SGARefLink{sga3:xvii:lemma:4-2-2}{4.2.2}) is trivial.
\item[(vi bis)]
  \SGARefTarget{sga3:xvii:proposition:4:3:1:item:vi-bis}
  \(\mathfrak g\) is nilpotent, and every subgroup of
  multiplicative type of the identity component of the center of
  \(G\) is trivial.
\item[(vi ter)]
  \SGARefTarget{sga3:xvii:proposition:4:3:1:item:vi-ter}
  \(G\) is nilpotent, and every subgroup of multiplicative type in
  the identity component of the center of \(G\) is trivial.
\end{enumerate}}
\SGARefTarget{sga3:xvii:proposition:4:3:1}

\SGAProof{Proof}
The implications
\SGARefLink{sga3:xvii:proposition:4:3:1:item:iii}{(iii)}\(\Rightarrow\)\SGARefLink{sga3:xvii:proposition:4:3:1:item:ii}{(ii)}\(\Rightarrow\)\SGARefLink{sga3:xvii:proposition:4:3:1:item:i}{(i)} are clear.  To prove
\SGARefLink{sga3:xvii:proposition:4:3:1:item:i}{(i)}\(\Rightarrow\)\SGARefLink{sga3:xvii:proposition:4:3:1:item:iii}{(iii)}, we shall need the following lemma.

\SGAStatement{\SGAInlineRefTarget{sga3:xvii:lemma:4-3-2}Lemma 4.3.2}{Let \(k\) be a field of characteristic
\(p>0\), let \(n\in\mathbb N\), and let \(k'\) be a radicial
extension of \(k\) such that \((k')^{p^n}\) is contained in \(k\).
For every \(k\)-prescheme \(X\) (respectively, every
\(k'\)-prescheme \(X'\)), denote by \(X^{(p^n)}\) (respectively,
\((X')^\phi\)) the \(k\)-prescheme obtained from \(X\)
(respectively, from \(X'\)) by the base change
\[
  F^n:k\longrightarrow k,\qquad x\longmapsto x^{p^n},
  \qquad
  \bigl(\text{respectively, }\;
  \phi:k'\longrightarrow k,\qquad x'\longmapsto(x')^{p^n}\bigr).
\]
Then, for every \(k\)-prescheme \(X\), there is a functorial
isomorphism
\[
  (X_{k'})^\phi\xrightarrow{\sim}X^{(p^n)}.
\]
Consequently, if \(X\) and \(Y\) are two \(k\)-preschemes for
which there exists a \(k'\)-isomorphism
\[
  u':X_{k'}\xrightarrow{\sim}Y_{k'},
\]
then there exists a \(k\)-isomorphism
\[
  v:X^{(p^n)}\xrightarrow{\sim}Y^{(p^n)}.
\]
If, moreover, \(X\) and \(Y\) are endowed with structures of
\(k\)-group preschemes and \(u'\) is a \(k'\)-homomorphism, then
\(v\) is a \(k\)-homomorphism.}
\SGARefTarget{sga3:xvii:lemma:4:3:2}

The lemma follows immediately from the transitivity of base
changes and from the fact that the composite
\[
  k\longrightarrow k'\xrightarrow{\phi}k
\]
is \(F^n\).
